\section{Discussion}
\label{sec:discussion}

\subsection{Lessons for engineering architectural digital twins}

\begin{itemize}[leftmargin=*,nosep]
  \item \textbf{L1: Generate the physical stand-in; do not build it by hand.} Generating the counterpart from $\mathcal{G}$ ensures traceability (\S\ref{sec:tracelinks}), co-evolution, and enumerable shared assumptions for Table~\ref{tab:independence}.
  \item \textbf{L2: Share one trace schema across simulation and observation.} Writing calibration and fidelity once for both sides eliminates dual-parser drift; it is cheapest early and costliest to retrofit.
  \item \textbf{L3: Fidelity evidence requires an independence analysis.} Table~\ref{tab:independence} prevents mistaking shared-plan agreement (e.g.\ exact topic rates) for validation; the informative findings lie strictly in independent rows (stall sets, liveliness, partition rediscovery, latency).
  \item \textbf{L4: Calibrators must abstain and disclose unidentifiable parameters.} Table~\ref{tab:identifiability} keeps realistic-scale evidence separate from unit round trips and prevents presenting unobservable staleness windows or periodic burstiness as calibrated.
  \item \textbf{L5: Make consistency debt visible.} One model instance is not enough; executable CI parity checks catch silent re-derivation bugs while an explicit debt register documents unfixed divergences.
  \item \textbf{L6: Match each service to its rung on the cost ladder.} Sub-second gates per commit, queue-flow screening, $\approx$48--54\,min/probe Docker \texttt{hybrid-verify} for milestone triage, and machine-hour counterpart sweeps for periodic twin validation keep the twin viable in CI/CD.
\end{itemize}

\subsection{From digital shadow to live digital twin (R7)}
\label{sec:livetwin}

The twin reaches the batch-shadow level (Table~\ref{tab:maturity}), with batch parameter/behavioural drift detection (\texttt{detect\_drift}, \texttt{diagnose\_fault\_residuals}) and incremental EMA re-calibration (\texttt{update\_profile}) evaluated on perturbed Docker counterparts (\S\ref{sec:evaluation}). That evaluation showed that rate/delivery drift and latency drift must be thresholded separately because single-machine latency can exceed 25\% even when delivery agrees within 0.3\%. Three live-twin capabilities remain \emph{designed} (stated in requirements voice):
\begin{itemize}[leftmargin=*,nosep]
  \item \textbf{Continuous ingestion:} The twin shall ingest OTel-shaped telemetry continuously, tagged with project, platform and version.
  \item \textbf{Topology drift:} Observed (application, topic) relations differing from $\Gs$ shall produce a graph diff and an engineer-reviewed \emph{model-repair proposal}, never an automatic rewrite.
  \item \textbf{Engineer-in-the-loop actuation:} The twin shall push accepted re-placements or configuration changes only through engineers---never autonomous production actuation, which is inappropriate for mission-critical middleware.
\end{itemize}
Building calibration (R5) and measured fidelity with known noise (R6) before live drift detection (R7) is forced: drift is a fidelity breach on new data.

\subsection{Threats to validity}
\label{sec:threats}

\paragraph{Construct validity.} $\Iobs$ adopts $\Idyn$'s delivery-loss definition rather than latency; stall-set Jaccard ignores partial loss; and only $\Idyn$ (not $I^*$, $\Icomp$ or $I_M$) was compared against the counterpart.

\paragraph{Internal validity.} Shared plan parameters, trigger functions and late-sample rules are bounded by Table~\ref{tab:independence}. Single-machine Python emulation adds scheduling noise, guarded by the capacity gate, CPU logging and WaitSet polling, though large topologies like \texttt{enterprise} hit single-machine discovery limits (\S\ref{sec:protocol}).

\paragraph{External validity.} All counterparts use generated timing stubs over real Cyclone DDS instead of a physical production system; broker semantics are relay-emulated, Kafka/MQTT/ROS~2 back-ends are not built, and real application business logic is untested. While this is an intentional, explicitly bounded pre-deployment trade-off (\S\ref{sec:independence}), validating against an independently written DDS or ROS~2 application under the same fault vocabulary is the key next external step.

\paragraph{Conclusion validity.} In Study~1 ($n = 8$ after blinded Amendment~2), exact Wilcoxon floors prevented H2 from clearing Holm ($p = 0.016 > 0.0125$) and allowed noisy \texttt{hub\_and\_spoke} ($C_s = 0.66$) to decide H1; moreover, O-open publisher crashes exercised mainly the dependent rows of Table~\ref{tab:independence}, explaining the tie with closed-form R. Study~2 ($n = 13$--$14$) addressed both limitations directly on independent rows, reaching confirmatory significance on steady-state latency (H2.1), infrastructure crashes and stall sets (H2.2a/b) and perturbed calibration (H2.3), while gray-failure omission error (H2.4) remained non-significant.
